\subsection{Hướng phát triển}
\label{sec:future_work}

Các hướng dưới đây xuất phát từ giới hạn quan sát được trong thực nghiệm, xếp
theo mức ảnh hưởng tới hệ thống hiện tại.

\begin{sloppypar}
\paragraph{Mã hóa ràng buộc đếm gọn hơn.}
Benchmark cho thấy chi phí giải tăng nhanh nhất theo kích thước $\mathit{KB}$,
và phần lớn mệnh đề đến từ mã hóa tổ hợp của \texttt{CountConstraint}, vốn sinh
$\binom{m}{k+1}$ mệnh đề cho một ràng buộc ``nhiều nhất $k$''. Các mã hóa dùng
biến phụ trợ như sequential counter \cite{sinz2005} cần $O(mk)$ mệnh đề, còn
totalizer \cite{bailleux2003} và sorting network \cite{batcher1968} cần
$O(m \log^2 m)$ đến $O(m^2)$ mệnh đề, đồng thời vẫn giữ khả năng lan truyền đơn
vị. Đây là thay đổi có lợi nhất nếu mở rộng lên lưới $6\times6$ trở lên.
\end{sloppypar}

\paragraph{Bộ giải CDCL và giải tăng dần.}
Bộ giải hiện tại quay lui theo thứ tự thời gian và dựng lại $\mathit{KB}_t$ từ
đầu ở mỗi truy vấn. Thuật toán CDCL \cite{marques1999grasp} học mệnh đề từ xung
đột và quay lui không theo thứ tự thời gian, giảm số lần phân nhánh trên các
bản đồ có ràng buộc toàn cục như \texttt{ConnectedConstraint}. Vì
$\mathit{KB}_{t+1}$ chỉ thêm mệnh đề vào $\mathit{KB}_t$
(Mục~\ref{sec:reveal_protocol}), một bộ giải tăng dần có thể giữ lại mệnh đề
đã mã hóa và mệnh đề đã học giữa các truy vấn thay vì bắt đầu lại.

\paragraph{Sinh bản đồ tự động.}
Sáu bản đồ hiện được thiết kế thủ công. Một bộ sinh tự động có thể chọn ngẫu
nhiên lời giải, sinh lời khai đúng với lời giải cho từng ô, rồi dùng
\texttt{validate\_puzzle} để giữ lại bản đồ có nghiệm duy nhất và không cần
đoán. Số sóng suy luận và tập dạng ràng buộc được phép dùng là hai tham số tự
nhiên để điều khiển độ khó.

\paragraph{Giải thích suy luận bằng ngôn ngữ tự nhiên.}
\texttt{DeductiveAgent} đã ghi lại, cho mỗi bước, các truy vấn SAT và lời khai
đang có hiệu lực. Bước tiếp theo là rút ra tập lời khai tối thiểu đủ để bác bỏ
trạng thái ngược lại, chẳng hạn bằng lõi không thỏa được (unsatisfiable core),
rồi diễn đạt nó thành câu giải thích cho người chơi.

\paragraph{Mở rộng dạng ràng buộc.}
Các dạng lời khai khác của trò chơi gốc, như ``người ở giữa hai người'', ``hàng
xóm chung'' hay ``nhiều nhất'' \cite{griductive_rules}, đều quy được về các
ràng buộc đếm trên một tập đích tính từ bàn chơi. Mỗi dạng mới chỉ cần một lớp
trong \texttt{constraints.py} và một hàm mã hóa trong \texttt{CNFEncoder}, không
cần thay đổi Game Engine.
